\documentclass[DIV=15]{scrartcl}

\usepackage{listings}
\usepackage{caption}
\usepackage{subcaption}
\usepackage{minted}
\usepackage{todonotes}
\usepackage{mathpartir}
\usepackage{amssymb}


\usepackage{caption}
\usepackage{subcaption}

\setminted{
  breaklines,
  breakanywhere,
  tabsize=2
}



\renewcommand{\`}[1]{\texttt{#1}}
\newcommand{\ket}[1]{|#1\rangle}
\newcommand{\code}[1]{\([[#1]]\)}


\title{A Verified Compilation Framework for Quantum Error-Corrected Systems}

\author{Andrés Paz, anpaz@cs.washington.edu}
%\author{Dan Grossman}
%\email{djg@cs.washington.edu}


% \begin{abstract}
% \end{abstract}

\begin{document}

\maketitle

% Include all sections
\section{Introduction}

Quantum computing promises exponential speedups for certain problems by exploiting superposition and entanglement~\cite{shor1997,wecker-reactions,cao2018potential,gidney2025factor2048bitrsa,resourceestimation-beverland2022assessing}. Yet writing reliable quantum software remains deeply challenging. Quantum programs are difficult to debug because each instruction updates probability amplitudes across all possible outcomes, and any measurement collapses the state, destroying the very information needed for inspection. 
Testing is equally limited: most quantum algorithms produce probabilistic results, requiring repeated runs to build statistical confidence. Because outputs are inherently nondeterministic and the input space grows exponentially, testing provides only limited assurance even in classical systems, and quantum computation simply amplifies this challenge. As a result, formal verification becomes essential for establishing correctness.


Quantum hardware itself adds further complexity. Different technologies (such as trapped ions~\cite{DeCross_2025}, superconducting circuits~\cite{google2025quantum}, majoranas~\cite{microsoft2025interferometric}, and neutral atoms~\cite{bluvstein2024logical}) realize quantum computation in distinct physical ways. Each provides its own native instruction set, reflecting the specific gate primitives supported by the hardware. The physical layout also matters: only certain pairs of qubits can directly interact, constraining how entangling operations are scheduled. Finally, all technologies are inherently noisy, and the types and rates of faults vary widely across implementations. As in classical computing, managing this hardware heterogeneity and noise requires systematic abstraction.

A standard or canonical instruction set abstracts away device-specific details, enabling programs to target multiple hardware platforms. Developers often assume idealized conditions such as all-to-all qubit connectivity, leaving the mapping to physical layouts for later stages. Quantum error correction (QEC) further strengthens this abstraction by encoding data across patches of many qubits, making logical operations more resilient to noise. From a programmer's perspective, these abstractions allow them to use perfectly behaved standard gates with all-to-all connectivity.

Yet these abstractions introduce a deep compilation challenge. Each transformation can radically alter a program's structure: for instance, a single logical CNOT on error-corrected qubits may expand into hundreds of physical gates with intermediate measurements and classical feedback. Classical control compounds this difficulty, as even simple protocols like teleportation require measurement-based branching, and error correction depends on real-time feedback from classical systems. A realistic compiler must therefore reason about both quantum and classical data across multiple abstraction layers, making correctness difficult to establish.

This complexity makes quantum compilers difficult to build and harder to trust. A single faulty transformation can silently alter program meaning, producing errors indistinguishable from hardware noise. Testing is insufficient, as outcomes are probabilistic and nondeterministic. Formal verification offers stronger guarantees but remains rare in practice. The verified compiler VOQC~\cite{Hietala_2021} embeds quantum semantics within Coq, ensuring correctness but requiring deep expertise and scaling poorly due to matrix-based reasoning over exponentially large states. More scalable approaches such as Giallar~\cite{tao2022giallar} use rewrite rules to check circuit equivalence, but these break down for programs involving measurements, feedback, or multiple abstraction levels. Moreover, most existing compilers treat QEC as a special-purpose layer rather than a first-class abstraction within the compilation pipeline, preventing correctness proofs from spanning logical and physical representations.

These challenges seem to demand a choice: either verify correctness rigorously at exponential cost, or scale efficiently without formal guarantees. However, this dilemma rests on an assumption, that verification must reason about quantum states directly. The key insight of this work is that correctness can instead be expressed over circuit structure. By defining formal semantics that describe how circuits transform logical information without expanding state vectors, we can reason about correctness efficiently while maintaining rigor. This enables proofs that compiler passes preserve meaning while avoiding exponential simulation costs, and it applies uniformly to hybrid quantum-classical programs and to QEC-protected circuits.

This thesis addresses these challenges by designing and implementing a compiler framework that integrates verification directly into the compilation process. Instead of verifying an entire compiler after construction, each compiler pass includes built-in checks that its transformations preserve program meaning. Verification thus becomes incremental and compositional, scaling naturally with program structure and extending across abstraction layers, including those introduced by QEC.

\paragraph{Thesis Statement.}
Verified compilation of quantum programs (including those with mid-circuit measurements, classical feedback, and quantum error correction) can be achieved compositionally by embedding structural correctness checks directly into each compiler pass, unifying formal assurance and practical compiler design across abstraction layers.

\vspace{2em}

This work makes the following contributions:

\begin{itemize}
  \item \textbf{A unified intermediate representation (IR).}  
  The IR supports multiple abstraction layers within the same program, treating error correction as a first-class hardware abstraction rather than a special-case extension. Programs can target physical or logical qubits interchangeably.

  \item \textbf{Classical–quantum interaction model.} The project defines a model that treats classical control as opaque callbacks, separating quantum semantics from classical computation while still supporting dynamic feedback. This separation simplifies reasoning about hybrid programs.

  \item \textbf{A compiler framework with built-in verification.}  Each compiler pass incorporates verification through forward-simulation and translation-validation techniques, integrating correctness directly into the compilation process rather than treating it as a separate validation step.
  
  \item \textbf{Formal semantics across abstraction layers.}  
  The IR defines compositional semantics that enable local proofs of correctness, ensuring transformations preserve meaning without full simulation.
  
  \item \textbf{Scalable verification techniques.}  
  Structural reasoning replaces matrix-based analysis, supporting programs with mid-circuit measurements, hybrid control, and QEC across multiple qubits.
  
  \item \textbf{Implementation and evaluation.}  
  A prototype compiler demonstrates verified compilation through multiple abstraction layers, from standard circuits to hardware-level implementations.
\end{itemize}

By embedding verification directly into compilation, this work unifies formal assurance and practical compiler design. It provides a foundation for scalable, architecture-aware quantum compilers that remain trustworthy as quantum programs and hardware systems grow in complexity.

The remainder of this proposal is organized as follows: Section~2 surveys related work in quantum languages, intermediate representations, and compiler verification. Section~3 presents the qstack framework and its core design principles. Section~4 describes the proposed work for this thesis. Section~5 explains how the work will be evaluated and concludes with a discussion of the scope of the work and future directions.




\section{Related Work}

This section surveys related work organized into three categories:

\begin{itemize}
\item \textbf{Languages and Intermediate Representations:} A survey of different languages and intermediate representations used for quantum programs, examining how they handle abstraction layers and classical control flow.
\item \textbf{Verification:} A survey of verification techniques for quantum programs and compilers, focusing on existing methodologies for checking correctness and their scalability limitations.
\item \textbf{Compilers:} A survey of quantum program compilers, with particular focus on those that handle quantum error correction codes and multi-layered compilation.
\end{itemize}

\subsection{Quantum Languages and Intermediate Represenations}


Research on quantum programming languages has evolved along two main axes: 
(1) how much abstraction a language provides to programmers, and 
(2) how it manages the boundary between classical and quantum computation. 
Across these axes, we can trace a progression from early circuit-construction frameworks, through statically typed, high-level imperative languages, to modern attempts that blend classical and quantum semantics within a unified model. 
This section summarizes representative approaches illustrating these trajectories.

One of the earliest and most influential high-level languages is \textbf{Quipper}~\cite{green2013quipper}. 
Implemented as an embedded domain-specific language in \textbf{Haskell}, Quipper allows programmers to describe quantum computations using familiar functional constructs while relying on the host language for classical control and meta-programming. 
A Quipper program is not directly executed on a quantum device: it is a \emph{circuit generator} that, when evaluated by the Haskell interpreter, produces a quantum circuit data structure. 
This circuit can then be optimized, serialized, or executed on a backend simulator or device.

Quipper’s core abstractions are minimal:
\begin{itemize}
  \item \textbf{Data types:} classical bits, qubits, gates, and boxed subcircuits.
  \item \textbf{Monad for circuit construction:} the \texttt{Circ} monad represented circuit-building computations.
  \item \textbf{Control:} a single quantum control construct, the \texttt{controlled} keyword, used to apply a gate conditionally on a classical or quantum bit. During execution, only this gate-level control is available; all other control flow occurs during circuit generation.
\end{itemize}

This design reflects the \emph{QRAM model}, where a classical controller orchestrates quantum operations. 
Its simplicity and explicit separation between classical meta-programming and quantum circuit generation made Quipper an ideal platform for expressing large-scale algorithms such as Quantum Fourier Transform, Shor’s algorithm, and simulation of physical systems. 
Quipper demonstrates that high-level, host-language embedding could scale to circuits with millions of gates, even though its execution model remained fully classical outside of gate control.

Building on the experience of circuit-generation frameworks, Q\#~\cite{qsharp} introduced one of the first stand-alone quantum programming languages with a full compilation and execution model. Earlier stand-alone efforts such as QCL~\cite{omer2005classical} and Scaffold~\cite{ScaffCC} explored similar directions, but Q\# was the first to support automatic generation of adjoint and controlled variants together with structured lifetime management. It also introduced practical tooling, including diagnostic primitives like \texttt{DumpMachine} and a lightweight unit-test framework for validating small programs.

Q\# enforces a strict separation between classical and quantum computation through its type system while supporting practical abstractions such as:
\begin{itemize}
  \item \textbf{Quantum operations and classical functions:} functions are purely classical and deterministic, while operations may act on qubits and produce quantum effects.
  \item \textbf{Automatic generation of adjoint and controlled variants:} developers only specify the forward definition, and the compiler automatically derives both adjoint and controlled versions when possible.
  \item \textbf{Scoped qubit lifetime management:} qubits are allocated with \texttt{use} or \texttt{borrow}, ensuring they are released automatically when their lifetime ends.
  \item \textbf{Structured uncomputation:} the \texttt{within}/\texttt{apply} pattern provides explicit uncomputation scopes, ensuring reversibility and preventing accidental measurement or loss of coherence.
  \item \textbf{Classical control flow:} loops and conditionals operate only on classical data, preserving a clean boundary between classical decision logic and quantum evolution.
\end{itemize}

This design defines a disciplined hybrid computation model: classical code manages orchestration and control, while quantum operations define unitary transformations. 
Its semantics remain close to the QRAM architecture but within a statically verified, first-class language.

As high-level quantum programming matured, new languages began exploring richer quantum abstractions that attempt to blur the classical–quantum boundary. 
\textbf{Silq}~\cite{silq} introduced \emph{safe, automatic uncomputation}---a mechanism allowing temporary quantum variables to be implicitly “cleaned up” without explicit programmer intervention. 
Its type system tracks which values remain entangled and ensures reversibility of operations without manual uncompute calls. 
Silq’s design philosophy sought to make quantum programming feel “classical” by unifying the treatment of classical and quantum data, but this unification also leaked into the syntax, introducing constructs like \texttt{!N} and \texttt{N} for classical and quantum versions of a variable. 
Rather than simplifying reasoning, these distinctions made programs harder to read and reason about. 
Although the semantics enforce safety, the resulting type system is intricate, combining linearity, ownership, and entanglement tracking.

Silq’s approach shows that many reversible computation patterns (such as the cleanup of ancillae in algorithms like Bernstein–Vazirani) can be handled automatically, but also illustrates the limits of abstraction: as one hides the quantum–classical distinction, the underlying computational effects become more opaque.


\textbf{Qunity}~\cite{qunity:Voichick_2023} gives quantum control first-class status.  
Control flow can depend on quantum data without measurement, so branches may execute coherently in superposition.  
A small core calculus enforces linearity and reversibility, and its type system ensures that coherent control is well defined.  
General recursion and function definitions are omitted because it conflicts with linearity and reversibility: non-terminating or non-unitary functions cannot be represented within a unitary semantics.  

To regain expressiveness, Qunity introduced a \emph{surface meta-language} built on top of the core\cite{compiling_qunity_2025}.  
This meta-language adds recursion, user-defined data types, and higher-level abstractions, but it is evaluated classically and operates outside the coherent fragment.  
Its role is to generate, compose, or structurally manipulate Qunity programs rather than to execute them coherently.  
This separation keeps the core reversible and analyzable while still supporting practical programming constructs through meta-level evaluation.

Qunity’s type system allows the compiler to statically determine where operations remain reversible and where measurement—and therefore classical control—is required.  
However, as in Silq, this typing discipline leaks into the syntax itself.  For example, Qunity provides several pattern-matching forms that differ in how they treat coherence:
\begin{itemize}
  \item \texttt{ctrl}: coherent control over possibly irreversible computations, allowed only when branches act on orthogonal subspaces and satisfy erasure conditions that prevent residual entanglement; preserves superposition.
  \item \texttt{pmatch}: pure, symmetric pattern match over orthogonal patterns; reversible and coherence-preserving.
  \item \texttt{ematch}: erasing match that preserves coherence when erasability conditions hold (the erased data is provably recoverable or irrelevant); otherwise disallowed in coherent contexts.
  \item \texttt{match}: classical match that may discard data and collapses superposition.
\end{itemize}
These constructs make the boundary between classical and coherent control explicit at each branch point.

While this explicitness preserves soundness, it shifts the burden of reasoning about coherence to the developer, making programs more verbose and sometimes harder to read—a central usability tension in the design.

\vspace{2em}

As higher-level languages explored these abstractions, another line of work focused on \emph{intermediate representations} that capture quantum and classical interactions close to hardware. 

\textbf{OpenQASM~3}~\cite{openqasm3} extends a simple quantum assembly language (QASM) family of intermediate representations originally designed to describe quantum circuits at the gate level.  
The original QASM was an informal text format for listing gate operations used in textbooks like \cite{nielsen&chuang}, later refined into \emph{OpenQASM~2}\cite{openqasm2}, which introduced a simple grammar and minimal classical features, focusing entirely on circuit-level descriptions.  
\emph{OpenQASM~3} generalizes this model by incorporating richer classical computation, timing primitives, and pulse-level instructions, with the goal of supporting real-time interaction between classical control and quantum hardware.

Its core logical circuits syntax defines:
\begin{itemize}
  \item \textbf{Classical datatypes:} \texttt{bit} and \texttt{bool}, together with new numeric types (\texttt{int}, \texttt{float}, \texttt{angle}, and \texttt{complex}) introduced in version~3 to enable real-time arithmetic and parameterized control.
  \item \textbf{Quantum datatypes:} \texttt{qubit} and quantum registers.
  \item \textbf{Quantum gates:} built-in single- and multi-qubit gates (e.g., \texttt{x}, \texttt{y}, \texttt{z}, \texttt{h}, \texttt{cx}, \texttt{rz}), measurement via \texttt{measure}, and reinitialization via \texttt{reset}.
  Starting with version~3, all gates are defined in terms of a single base gate \texttt{U}, whose semantics cover arbitrary single-qubit unitaries.  
  A standard library of canonical gates (e.g., \texttt{x}, \texttt{cx}, \texttt{rz}) is provided as built-in definitions over \texttt{U}.
\end{itemize}

Beyond pure circuit specification, OpenQASM~3 adds several constructs for classical and timing behavior:
\begin{itemize}
  \item \textbf{Classical expressions and assignments:} arithmetic, logical, and comparison operators restricted to a small, implementation-oriented subset. More general computation can be delegated to host code via \texttt{extern}, which implements near-time control by invoking classical functions between gates evaluation.
  \item \textbf{Control flow:} loops and conditional branches that may depend on measurement outcomes.
  \item \textbf{Timing primitives:} explicit scheduling with \texttt{delay}, \texttt{duration}, and timing annotations (\texttt{@t}) for device-level coordination.
  \item \textbf{Pulse-level primitives:} low-level operations that specify hardware control pulses, enabling descriptions of physical circuits rather than only logical ones.
\end{itemize}

While OpenQASM~3 is often described as a multi-level IR, this layering primarily spans logical (gate-level) and physical (pulse-level) program descriptions, rather than distinct quantum instruction sets.  
Its expanded control structures, numeric types, and timing model were largely shaped by IBM’s hardware and control architecture, reflecting pragmatic rather than purely language-theoretic design choices.  
The result is a unified but hardware-oriented language that bridges circuit compilation, hardware instructions, and real-time device control within a single syntactic framework.




As OpenQASM~3 defines a textual, multi-level language, the \textbf{Quantum Intermediate Representation (QIR)}~\cite{qir-spec} serves as its low-level, compiler-oriented counterpart. 
QIR is an extension of the LLVM Intermediate Representation designed to represent quantum programs in a form suitable for optimization and backend compilation. 
It provides a standardized way for quantum frontends to target diverse hardware backends through a common IR layer, allowing compilers, simulators, and runtime systems to interoperate.

QIR defines two main execution \emph{profiles}:
\begin{itemize}
  \item The \textbf{Base profile} restricts programs to straight-line quantum code ending in measurement, suitable for precompiled circuits without feedback.
  \item The \textbf{Adaptive profile} extends this to allow mid-circuit measurements, classical branching, and data-dependent feedback on qubits.
\end{itemize}
Profiles constrain the allowed LLVM instructions, control flow, and runtime intrinsics to ensure compatibility with specific device capabilities. 
The design separates platform-independent quantum semantics from hardware-specific implementations, letting compilers generate profile-compliant QIR modules and letting backends interpret or translate them into native instruction sets. 
Extensions such as \texttt{Reset}, \texttt{Integer}, and \texttt{FloatingPoint} add incremental support for mixed classical–quantum computation.

\vspace{2em}

While languages like OpenQASM~3 and QIR aim to generalize across hardware and execution models, \textbf{Stim}~\cite{gidney2021stim} targets a narrower but crucial domain: efficient simulation of stabilizer circuits used in quantum error correction.
Stim is designed not as a general programming language but as a \emph{fast stabilizer-circuit simulator}. 
Stim focuses on a subset of quantum computation—\emph{Stabilizer circuits}, consisting of Cliffords gates with measurements and classically controlled Pauli corrections—which can be simulated efficiently on classical hardware.

Clifford circuits are built from gates in the \{H, S, CNOT\} set, which preserve the stabilizer structure of quantum states. 
They are not universal for quantum computation, since they cannot generate arbitrary unitary transformations (most notably, the T gate is missing), but they form the foundation of most quantum error correction protocols. 
Because they admit an efficient classical simulation through tableau representations, Clifford circuits strike a balance between expressivity and scalability.

Stim’s language is intentionally minimalist:
\begin{itemize}
  \item No general variables or functions beyond the circuit structure itself.
  \item Measurements produce classical outcomes that can control later gates, written as conditionals like \texttt{CX rec[-1] 1}, where \texttt{rec[-1]} references the most recent measurement result.
  \item A \texttt{REPEAT} instruction allows replication of subcircuits, enabling loops without introducing dynamic control flow.
\end{itemize}

Because of this simplicity, Stim achieves extraordinary performance. 
According to its authors, it can simulate a surface-code circuit with 20{,}000 qubits, 8 million gates, and 1 million measurements in about 15 seconds, and generate measurement samples at roughly 1~kHz. 
These capabilities have made it the \emph{de facto} standard simulator for quantum error correction research. 
Its efficiency stems from a carefully engineered tableau representation and from limiting expressivity to the stabilizer subset, which admits polynomial-time simulation. 
Stim demonstrates the other extreme of the design spectrum: rather than raising abstraction, it narrows the domain to achieve scalability, providing essential infrastructure for verification, decoding, and fault-tolerance analysis in large-scale quantum error correction.

\begin{figure}[p]
\centering
\newcommand{\lcol}{0.52\linewidth}
\newcommand{\rcol}{0.44\linewidth}

% ---------- Row 1 ----------
\begin{subfigure}[t]{\lcol}
\centering
\textbf{Quipper}
\begin{minted}{haskell}
bell :: Circ (Qubit,Qubit)
bell = do
  a <- qinit False; b <- qinit False
  hadamard a
  qnot_at b `controlled` a
  return (a,b)

with_computed bell $ \(a,b) -> do
  qnot_at a `controlled` b
  cdiscard (a,b)
\end{minted}
\end{subfigure}\hfill
\begin{subfigure}[t]{\rcol}
\centering
\textbf{Qunity}
\begin{minted}{text}
-- ASCII-only surface style
def deutsch(f : Bit -> Bit) : Bit =
  let x = plus in
  ctrl f(x) [
    0 -> x ;
    1 -> { gphase(pi); }
  ];
  had x;
  return x;
end
\end{minted}
\end{subfigure}

\vspace{2em}

% ---------- Row 2 (left only) ----------
\begin{subfigure}[t]{\linewidth}
\raggedright
\textbf{Q\#}
\begin{minted}{csharp}
operation CRy(theta : Double, ctrl : Qubit, tgt : Qubit)
  : Unit is Adj + Ctl {
    (Controlled Ry)([ctrl], (theta, tgt));
}

operation PrepareBell() : Unit {
    use (q0, q1) = (Qubit(), Qubit());
    H(q0);
    CNOT(q0, q1);
}
\end{minted}
\end{subfigure}

\vspace{2em}

% ---------- Row 3 ----------
\begin{subfigure}[t]{\lcol}
\centering
\textbf{OpenQASM 3}
\begin{minted}{c}
OPENQASM 3.0;
qubit q[2];

gate x(a) { U(pi, 0, pi) a; }
gate h(a) { U(pi/2, 0, pi) a; gphase(-pi/4); }

h q[0];
ctrl @ x q[0], q[1];   // CNOT via modifier
\end{minted}

\end{subfigure}\hfill
\begin{subfigure}[t]{\rcol}
\centering\textbf{Stim}
\begin{minted}{text}
H 0
H 1
MPP X0*X1
DETECTOR rec[-1]
OBSERVABLE_INCLUDE(0) rec[-1]
\end{minted}
\end{subfigure}

\vspace{2em}


% ---------- Full-width bottom ----------
\begin{subfigure}[t]{\linewidth}
\raggedright
\textbf{QIR}
\begin{minted}{llvm}
%q0 = call %Qubit* @__quantum__rt__qubit_allocate()
%q1 = call %Qubit* @__quantum__rt__qubit_allocate()

call void @__quantum__qis__h__body(%Qubit* %q0)
call void @__quantum__qis__cnot__body(%Qubit* %q0, %Qubit* %q1)

call void @__quantum__rt__qubit_release(%Qubit* %q1)
call void @__quantum__rt__qubit_release(%Qubit* %q0)
\end{minted}
\end{subfigure}

\caption{
Six minimal snippets foregrounding each approach to representing circuits:
(a) Quipper’s boxed subcircuits and controlled combinators;
(b) Qunity’s coherent control over quantum data;
(c) Q\# functors for controlled/adjoint variants;
(d) OpenQASM~3 gate definitions and modifiers;
(e) Stim’s stabilizer format with parity measurements and detectors;
(f) QIR’s LLVM calls to quantum intrinsics.
}
\label{fig:lang-contrast-6}
\end{figure}



\subsection{Verification of Quantum Programs and Compilers}
\label{sec:verification}

Verifying the correctness of quantum programs is essential as circuits grow in complexity and manual reasoning or testing becomes infeasible. 
Testing, the standard practice in classical software engineering, fails in the quantum setting for both theoretical and practical reasons. 
As Rovara et al.\ show in their framework for debugging quantum programs~\cite{rovara2025framework}, even identifying that an error occurred is difficult, let alone determining its cause. 
The authors argue that debugging quantum software is fundamentally more difficult than classical debugging for several reasons:

\begin{enumerate}
  \item A single quantum operation affects the entire system state due to entanglement.
  \item Measurement collapses superposition, making it impossible to inspect intermediate states nondestructively.
  \item The global state space grows exponentially with the number of qubits, making full simulation intractable.
\end{enumerate}

Their system introduces assertion-like constructs that can check limited properties of quantum states, such as orthogonality or separability, but these checks require access to the full quantum state through simulation. As a result, they are only practical for very small programs and cannot be used to test or debug realistic quantum circuits. This reflects a broader limitation of testing itself: even in classical software, nondeterminism and the combinatorial explosion of inputs prevent exhaustive validation—quantum programs simply magnify the problem. Because quantum states cannot be inspected without measurement, and measurement destroys the information being observed, formal reasoning remains the only reliable path to establishing correctness.


\vspace{1em}

One of the first steps toward proving correctness in the quantum setting came with \textbf{Quantum Hoare Logic (QHL)}~\cite{ying2019hoare}. 
QHL extends classical Hoare logic\cite{hoare1969axiomatic}  to the quantum domain, introducing preconditions and postconditions expressed as Hermitian operators rather than Boolean predicates.
A Hoare triple of the form $\{A\}\,P\,\{B\}$ expresses that executing program~$P$ transforms an input state satisfying observable~$A$ into a state satisfying~$B$.
This allows reasoning about programs algebraically rather than through explicit simulation.
While elegant, early formulations of QHL were limited to closed quantum systems without measurement or classical control.
Later work\cite{feng2021quantum} extended the logic to handle partial measurements and mixed quantum–classical states, but scalability remained a challenge.
These logics laid the conceptual foundation for formal reasoning about quantum programs, but they stopped short of providing verified compilers that could connect high-level programs to low-level circuits.

\vspace{1em}

The first verified compiler in this space is \textbf{VOQC}~\cite{Hietala_2021}, which builds on these ideas to verify quantum circuit transformations in the proof assistant Coq. 
VOQC introduces two embedded languages, QWIRE and SQIR, that represent circuits at different abstraction levels and defines circuit semantics via their corresponding matrices. 
Assertions about circuit equivalence are then expressed as matrix equalities, allowing the use of formal proofs in Coq.
To show that compiler passes preserve semantics, VOQC adopts the same \emph{forward simulation} principle used in classical verified compilers such as \textbf{CompCert}~\cite{leroy2009compcert}. 
In forward simulation, one proves that each step of a compiled program corresponds to one or more steps of the source program that preserve observable behavior.
Given source and target semantics $\rightarrow_S$ and $\rightarrow_T$, we define a relation~$R$ such that:
\[
  s_S \;R\; s_T \;\wedge\; s_S \rightarrow_S s_S'
  \;\Rightarrow\;
  \exists s_T'.\; s_T \rightarrow_T^* s_T' \wedge s_S' \;R\; s_T' .
\]
If the source semantics is deterministic, then every terminating source execution has a matching target execution that produces an equivalent observable result.
This provides a modular and compositional proof structure: each compiler pass is verified independently, and correctness composes across passes.
VOQC thus establishes the first fully verified quantum optimizing compiler, directly extending the CompCert tradition into the quantum domain.

Despite its strong formal guarantees, VOQC faces two major limitations.
First, because equivalence is expressed as equality between $2^n \times 2^n$ matrices, verification scales poorly beyond a handful of qubits.
Second, its deep embedding in Coq makes it cumbersome to develop and experiment with real-world circuits or compiler passes.
These challenges motivated more pragmatic approaches to verification.

\vspace{1em}

\textbf{Giallar}~\cite{tao2022giallar} takes a pragmatic turn by targeting the \texttt{Qiskit} compiler, one of the most widely used quantum programming environments.
Rather than defining a new language or embedding circuits in a proof assistant, Giallar verifies the correctness of existing compiler passes directly.
Its verification model is structured around two complementary ideas:

\begin{enumerate}
  \item \textbf{Verified rewrite rules:} it define circuit rewrite rules (e.g., gate commutation, decomposition, or cancellation) which are individually proved correct once and for all in Coq.
  \item \textbf{Pass verification:} a compiler pass is correct if its transformation can be expressed entirely as a sequence of these verified rewrites.
\end{enumerate}

This approach reduces the problem of verifying large compiler passes to checking whether their transformations decompose into trusted rewrite steps.
If the input circuit can be transformed into the output circuit through these rules, equivalence follows immediately.
The method scales far better than matrix-based equivalence checking because it relies on syntactic and structural reasoning rather than full state representations.
Beyond verifying quantum transformations, Giallar also analyzes the \emph{classical} control logic of compiler passes.
By symbolically executing control flow and constraints in combination with the Z3 solver, it detects mismatches between the intended and actual behavior of compiler passes.
In their evaluation, the authors verified $44$ out of $56$ passes in \texttt{Qiskit}, uncovering several previously unknown bugs.
Unsupported cases include randomized or heuristic optimizations, approximate synthesis, and pulse-level transformations, which lie outside the formal model.

\vspace{1em}

\vspace{1em}

All the approaches described so far reason in the Schrödinger picture, where programs are defined by how they transform quantum \emph{states}. 
This representation becomes costly as circuit size grows, since each $n$-qubit state requires tracking $2^n$ amplitudes. 
An alternative is to reason in the \emph{Heisenberg picture}, where programs act on \emph{observables} instead. 
This shift replaces state evolution with operator transformation and forms the basis of \textbf{Hoare Meets Heisenberg: A Lightweight Logic for Quantum Programs}~\cite{sundaram2022hoare}.

Formally, in the Schrödinger picture applying a unitary~$U$ to a quantum state~$\rho$ transforms it as
\[
  \rho' = U\rho U^\dagger,
\]
while in the Heisenberg picture, an observable~$O$ evolves as
\[
  O' = U^\dagger O U.
\]
Because $\mathrm{Tr}(\rho' O) = \mathrm{Tr}(\rho O')$, both perspectives yield the same measurement statistics, but the Heisenberg view allows reasoning over observables instead of exponentially large states.

This shift makes reasoning about programs symbolic rather than numeric. 
Instead of simulating how a quantum state evolves step by step, the logic tracks how each operation transforms observables and postconditions. 
Correctness is established by checking that the given precondition~$A$ implies the weakest precondition obtained by propagating the postcondition~$B$ backward through the inverse of each gate.

For example, consider verifying a Hadamard gate swaps the $Z$ and $X$ measurement bases, two Hoare triples are considered:
\[
\{Z\}\,H\,\{X\}
\qquad\text{and}\qquad
\{X\}\,H\,\{Z\}.
\]
Each specifies a precondition (the observable before the gate) and a postcondition (the observable after it). 
Verification proceeds by propagating each postcondition backward through the inverse of the gate. 
Since $H^\dagger X H = Z$ and $H^\dagger Z H = X$, both triples hold—confirming that the Hadamard exchanges the $Z$ and $X$ measurement bases as expected.


This backward reasoning avoids simulating state evolution altogether: it manipulates the symbolic form of measurement operators directly. 

Because Clifford gates map Pauli operators to other Paulis, the entire process reduces to algebraic manipulation of Pauli strings, which scales linearly with circuit size. 
This idea closely mirrors the approach used by stabilizer simulators such as \textbf{Stim}~\cite{gidney2021stim}, which track Pauli transformations rather than amplitude vectors to achieve large-scale efficiency.

\vspace{1em}

The same Heisenberg-based reasoning is further formalized and generalized in \textbf{Stabilizer Circuit Verification}~\cite{kliuchnikov2023stabilizer}.
This framework verifies equivalence between circuits—including those with measurement and classical feedback—by comparing how they transform stabilizer generators and classical outcomes, rather than quantum states.
A circuit’s behavior is captured by its \emph{logical action}, describing how logical operators of a code (such as $X_L$ or $Z_L$) are mapped to new ones.
For example, in a three-qubit repetition code, the logical operator $Z_L = Z_1Z_2Z_3$ represents the parity of all three qubits.
An encoding circuit might map $Z$ to $Z_1Z_2Z_3$, while a decoding circuit performs the inverse.
Two circuits are logically equivalent if they induce the same transformation on all stabilizers and measurement outcomes.

Under the mild restriction that classical control depends only on the \emph{parity} of prior measurements, these transformations can be computed efficiently using matrices over~$\mathbb{F}_2$, avoiding the exponential branching of arbitrary control.
This representation enables verification of circuits far beyond the reach of full-state simulation, including lattice-surgery operations, code encoders and decoders, and inter-code transformations.
Although motivated by quantum error correction, the method generalizes naturally: any circuit can be seen as acting on a logical code.
Even a single-qubit physical circuit implements the trivial $[[1,1,1]]$ code—one logical qubit encoded into one physical qubit.
From this perspective, every circuit performs a logical action, and verifying transformations means confirming that these actions are identical.
This unifies physical and logical verification: the same reasoning used to check code-switching or fault-tolerant operations also applies to compiler transformations and circuit optimizations.

\subsection{Compilers for Quantum Programs}


\textbf{Qiskit}~\cite{javadiabhari2024quantumcomputingqiskit} is often introduced as a Python-embedded domain-specific language for constructing quantum circuits, but this is only its surface layer. The language front end is primarily a wrapper for OpenQASM, while Qiskit’s main contribution lies in its modular compiler and execution stack. 

Its core abstractions include the following components:

\begin{enumerate}
  \item \textbf{Circuits:} the fundamental data structure representing quantum computations. Circuits can be serialized to formats such as OpenQASM or QIR.
  \item \textbf{Transpiler framework:} a retargetable system that converts circuits into device-specific forms. Compiler transformations are expressed as ordered sequences of \emph{passes}, grouped within \emph{pass managers}.
  \item \textbf{Passes and pass managers:} each pass can perform tasks such as high-level synthesis, qubit routing, or low-level optimizations like gate commutation and cancellation.
  \item \textbf{Target model:} encapsulates hardware-specific properties including the instruction set, qubit connectivity, timing constraints, and error rates. This model guides synthesis, mapping, and scheduling.
  \item \textbf{Execution primitives:} abstractions such as the \texttt{Sampler} and \texttt{Estimator} define a unified execution interface across simulators and physical devices, supporting parameter binding, expectation-value estimation, and built-in error mitigation.
\end{enumerate}


Qiskit also provides mechanisms for dynamic circuits, allowing classical conditions and mid-circuit measurements, and uses Rust-based implementations for performance-critical passes. Beyond its language surface, Qiskit functions as a flexible, pass-based compiler infrastructure that integrates analysis, optimization, and backend targeting.

\textbf{QVM}~\cite{qvm2024} builds on this infrastructure and demonstrates how new abstractions can extend circuit compilation. It introduces the idea of \emph{Gate Virtualization}, where selected two-qubit gates are replaced by sequences of single-qubit operations and classical post-processing. This allows large circuits to be split into smaller fragments that can run on limited hardware or across multiple devices.

To manage this, QVM defines a new intermediate representation called the \emph{Virtual Circuit IR (VC-IR)}. It augments Qiskit’s circuit model with virtual gates and explicit fragments, making boundaries and dependencies visible to the compiler. Dedicated passes handle circuit cutting, dependency reduction, and qubit reuse. Each fragment runs separately, and the runtime recombines results using classical correlations. The approach increases effective qubit capacity and maintains logical equivalence, showing how Qiskit’s pass system can support partitioned, hybrid computation.

\vspace{1em}

\textbf{Automated Synthesis of Fault-Tolerant State Preparation Circuits for Quantum Error-Correction Codes}~\cite{peham2025automated} is part of the \textbf{MQT-QECC} project~\cite{wille2024mqt}, a toolkit developed at Technical University Munich (TUM) for automating the design of quantum error-correction circuits. The paper focuses on how to automatically build reliable circuits that prepare encoded logical states. In quantum computing, \textif{fault tolerance} means that a circuit continues to work correctly even when some gates or qubits fail—specifically, a single physical error must not spread and cause multiple logical errors. The authors propose a synthesis method that enforces this property by design: it automatically generates both a preparation circuit and a verification step that together ensure any single fault is detected. The circuit repeats the preparation and verification until it passes all checks, guaranteeing correctness. This approach produces efficient circuits for common error-correcting codes and illustrates how compiler tools can directly enforce physical reliability, not just optimize performance.


\vspace{1em}

\textbf{A Game of Surface Codes}~\cite{litinski2019game} introduces a visual and structured way to describe computation using \textie{surface codes}, one of the most practical forms of quantum error correction. A surface code arranges qubits on a 2D grid where groups of qubits collectively store a single logical qubit, continuously detecting and correcting physical errors through stabilizer measurements. Because surface codes can tolerate high error rates and fit well with planar hardware layouts, they are considered a leading approach for scalable quantum computing.

Instead of describing programs as sequences of gates, the paper models computation as the manipulation of rectangular \emph{patches} on this grid. Logical operations are carried out by merging, splitting, and deforming these patches, a process known as \textit{lattice surgery}. Each operation corresponds to a simple measurement pattern, making the logic both spatial and temporal.

The paper also outlines how this geometric computation could be compiled in stages: first by mapping logical gates to surgery primitives, then by arranging them in space and time to respect hardware constraints. In doing so, it proposes a multi-pass view of compilation—moving from abstract operations to concrete, fault-tolerant layouts—and provides a clear blueprint for how future compilers can bridge the gap between logical algorithms and physical surface-code implementations.

\textbf{Lattice Surgery Compilation Beyond the Surface Code}~\cite{herzog2025latticesurgerycompilationsurface} turns that blueprint into an operational compiler. It generalizes lattice-surgery compilation to multiple topological codes through the notion of a \emph{code substrate}.
A substrate defines allowed surgery operations and connectivity constraints for a given code, represented as a routing graph. Compilation becomes the task of mapping a logical Clifford+T circuit onto this graph, generating valid measurement sequences and space–time layouts.
Unlike Litinski’s conceptual work, this paper delivers an implementation within the MQT framework. It takes standard logical circuits—often in OpenQASM or similar form—and emits JSON-based schedules and layouts.
While still limited to stabilizer-based computations, it demonstrates how geometric compilation can be automated and extended across code families. 


\section{\texttt{qstack}}
\label{sec:qstack}

\newcommand{\keval}{\Rightarrow }
\newcommand{\qeval}{\mapsto_q }
\newcommand{\ceval}{\mapsto_c }

\texttt{qstack} is a project developed to explore and validate many of the concepts that this thesis builds upon. 
It demonstrates that a quantum circuits language with opaque classical control is not only possible, but simplifies reasoning about programs across different abstraction levels. 
By defining computation purely in terms of quantum instructions, \texttt{qstack} provides a minimal yet expressive foundation with well-defined quantum semantics and a clean, abstract interface to classical control through opaque callbacks, independent of any hardware or instruction set.


\subsection{Language Model}

\texttt{qstack} is designed around the idea that every quantum computation can be expressed as a composition of simple \emph{quantum kernels}. 
Each kernel describes a self-contained computation that allocates a qubit, performs a sequence of quantum instructions, measures that qubit, and finally calls a classical function that determines what to do next. 
This cycle  of \texttt{allocate} $\rightarrow$ \texttt{compute} $\rightarrow$ \texttt{measure} $\rightarrow$ \texttt{callback}, defines a complete execution model that requires no classical expressions within the language itself. 
The result is a purely quantum core with exact and compositional semantics.

A \texttt{qstack} program consists of nested kernels. 
The general form of a kernel is:

\begin{minted}{ruby}
allocate q:
  instructions
measure callback
\end{minted}

Here, \texttt{allocate q} introduces a new qubit initialized to the $\ket{0}$ state. 
The body contains one or more quantum instructions or nested kernels. 
The kernel ends with a measurement of the corresponding allocated qubit and passes the classical measurement result to a callback function, implemented externally (for instance, in Python). 
By default, the keyword \texttt{measure} without arguments uses the callback \texttt{collect}, which stores measurement results in a classical stack. 

Because every kernel measures its allocated qubit, \texttt{qstack} programs follow a strict stack discipline for resource management: qubits are allocated and deallocated in last-in, first-out order. 
This keeps allocation, measurement, and control flow local, avoiding programmer-visible global state and any classical manipulation of values (though the semantics maintains an internal mapping of qubit identifiers to positions).

\subsection{Example}

A minimal program that demonstrates dynamic control flow is the repeat-until-success (RUS) pattern, where computation continues until a measurement yields a desired outcome:

\begin{minted}{ruby}
allocate q1:
  allocate q2:
    mix q2
    entangle q2 q1

  measure collect
measure repeat_until_zero
\end{minted}

The \texttt{repeat\_until\_zero} callback is a classical function that inspects the most recent measurement and decides whether to continue the computation. 
If the result is~1, it returns a new kernel identical to the original; otherwise, it let's the program to continue by returning \texttt{None} and the program terminates. 
This illustrates how probabilistic control flow can be expressed entirely within quantum semantics, with the callback mechanism acting as an opaque bridge to classical computation.


\subsection{Formal Semantics}

\texttt{qstack} defines a complete small-step operational semantics that specifies how programs evolve during execution. 
A program configuration is represented as $(C, S, (|\psi\rangle, m))$, where:
\begin{itemize}
    \item $C$ is the stack of active kernels and their associated callbacks,
    \item $S$ is the classical store (opaque to the semantics),
    \item $(|\psi\rangle, m)$ is the current quantum state vector and the mapping of logical qubit identifiers to their positions in the state.
\end{itemize}

Program execution proceeds as a sequence of transitions:
\[
C, S, (|\psi\rangle, m) \keval C', S', (|\psi'\rangle, m')
\]
where $\keval$ denotes one small-step of evaluation.

The semantics is defined by four inference rules, describing how kernels are started, how gates are applied, and how kernels complete execution with or without continuation.

\vspace{1em}

\noindent
\textbf{Start-kernel:} allocates a new qubit in the $\ket{0}$ state, extends the state vector, and begins execution of the kernel’s instruction list. 
\vspace{-1em}
\begin{center}
\begin{mathpar}
\inferrule[start-kernel]
{ }
{ ((q_k:I_k:cb_k): I, cb):C', S, (\ket{\psi^n}, m)  \keval (I_k, cb_k):(I, cb):C', S, (\ket{\psi^{n}} \otimes \ket{0},\ m[q_k \mapsto n])  }
\end{mathpar}
\end{center}

\noindent
\textbf{Quantum-gate:} applies the unitary operator associated with the current gate to its target qubits. 
The quantum state $\ket{\psi}$ is updated by applying the lifted unitary $U_{gate}$ to the indices indicated by the mapping $m$. 
\vspace{-1em}
\begin{center}
\begin{mathpar}
\inferrule[quantum-gate]
{  }
{ (\mathtt{gate\_id}(q_0, \ldots, q_j): I', cb):C', S, (\ket{\psi}, m) \keval (I', cb):C', S, (U_\mathtt{gate\_id}^{(m[q_0], \ldots, m[q_j])}\ket{\psi}, m)  }
\end{mathpar}
\end{center}



\noindent
\textbf{End-of-kernel-done:} handles the completion of a kernel when its instruction list is empty.  
The top qubit is measured, and the resulting classical bit is passed to the callback that returns no further kernel, execution proceeds by removing the completed kernel from the stack and continuing with the next element.
\vspace{-1em}
\begin{center}
\begin{mathpar}
\inferrule[end-of-kernel-done]
{  
    \|\psi^{n}\|_{n-1} \mapsto b, \ket{\psi'} \\ 
    m =   m' \uplus  [q \mapsto n-1]  \\
    f_{cb}(b,S) \mapsto \varnothing, S'  }
{ ([], cb):C', S, (\ket{\psi^n}, m) \keval C', S', (\ket{\psi'}, m') }
\end{mathpar}
\end{center}

\noindent
\textbf{End-of-kernel-continue:} also applies when a kernel has no remaining instructions.  
The top qubit is measured and the result is passed to the callback, which in this case returns a new kernel.  
That new kernel is immediately placed on top of the execution stack and evaluated, extending the quantum state with a fresh $\ket{0}$ qubit.
\vspace{-1em}
\begin{center}
\begin{mathpar}
\inferrule[end-of-kernel-continue]
{ \|\psi^{n}\|_{n-1} \mapsto b, \ket{\psi'} \\
    m =   m' \uplus  [q \mapsto n-1]  \\
    f_{cb}(b,S) \mapsto (q_k:I_k:cb_k), S'  }
{ ([], cb):C', S, (\ket{\psi^n}, m) \keval (I_k, cb_k):C', S',  (\ket{\psi'} \otimes \ket{0},\ m'[q_k \mapsto n-1])  }
\end{mathpar}
\end{center}


\smallskip
Together, these rules define the complete lifecycle of a kernel---allocation, computation, measurement, and continuation---capturing the precise operational behavior of \texttt{qstack} programs.


\subsection{Compositional Compilation and Error Correction}

A key result of the \texttt{qstack} project is that this kernel model naturally supports compositional compilation. 
Each compiler phase acts as a transformation on quantum programs: it may change the instruction set or insert additional operations, but the resulting program is still evaluated under the same operational semantics. 
That is, every compiled program remains a valid \texttt{qstack} program whose behavior is defined by the same small-step rules for allocation, gate application, measurement, and callback execution. 
This uniform semantic framework allows different compiler passes to be combined modularly, even if the transformations they apply are substantial.

Callbacks make this process systematic. 
Since callbacks are treated as opaque classical functions, the compiler does not inspect or modify their internal logic. 
When a transformation changes how measurements are represented—for example, when an error-correcting code expands a single logical measurement into several physical ones—the compiler automatically \emph{wraps} each callback with an adapter function. 
This wrapper decodes the measurement results into the form expected by the original callback and recompiles any kernels that the callback dynamically generates. 
Through this mechanism, callbacks remain compatible across compilation layers, and all resulting programs continue to execute according to the same formal semantics.

Using this infrastructure, \texttt{qstack} implements complete compilers for quantum error correction codes, including the three-qubit repetition code~\cite{nielsen&chuang} and the Steane [[7,1,3]] code~\cite{steane1996error}. 
The repetition compiler expands single-qubit operations into triple encodings with majority-vote decoding, while the Steane compiler injects stabilizer measurement routines, syndrome extraction, and decoding logic. 
Both integrate into an end-to-end compilation pipeline that begins with \texttt{qstack}'s toy instruction set, applies one or more levels of concatenated error correction, and finally targets hardware-specific instruction sets such as those used by trapped-ion devices. 
Each compilation step produces a new \texttt{qstack} program evaluated under the same semantics, enabling reasoning about compiler correctness within a single, coherent execution model.



\section{Proposed Work}

For this thesis I propose to build a \textbf{verified compiler framework} for quantum programs grounded in the \textbf{operational semantics of \texttt{qstack}}.  
The goal is to formally validate the correctness of both \textbf{instruction-level} and \textbf{whole-circuit} compiler passes without relying on interactive theorem provers.  
Instead, the project will use a combination of \textbf{stabilizer-based logical-action verification} and \textbf{density-matrix equivalence checking} to automatically establish semantic preservation.  
For scalability, we prefer the former, so we will develop a tool to determine whether a circuit belongs to the stabilizer formalism.

\subsection{Language Design}

We will define a \textbf{minimal quantum language} derived from OpenQASM~3 that preserves the same \textbf{quantum/classical separation} as \texttt{qstack} while removing features that complicate formal reasoning, such as dynamic allocation and timing control.  
The retained constructs include:
\begin{itemize}
    \item A single universal \textbf{base gate};
    \item \textbf{User-defined composite gates} built from existing ones;
    \item \textbf{Static classical and quantum registers};
    \item \textbf{Gate application}, \textbf{measurement}, and \textbf{reset};
    \item \textbf{Conditionals} based on XORs of classical bits; and
    \item \textbf{\texttt{extern}} declarations for opaque operations.
\end{itemize}

This subset remains expressive enough to encode all \texttt{qstack} kernels and the \textbf{stabilizer circuits} defined in~\cite{kliuchnikov2023stabilizer}.  
It provides a compact foundation for formal semantics and automated verification.

\subsection{Formal Semantics}

The formal semantics will be inspired on the \textbf{small-step operational semantics} of \texttt{qstack}.  
A program configuration consists of:
\begin{itemize}
    \item A \textbf{quantum state} represented as a density matrix $\rho$;
    \item A \textbf{mapping} from logical qubit identifiers to physical indices; and
    \item A \textbf{classical store} containing deterministic registers and callback state.
\end{itemize}

Gate applications correspond to \textbf{unitary superoperators}, measurements collapse the relevant qubit and update the store, and callbacks may determine continuation through classical control.  
This semantics supports reasoning about both low-level physical transformations and high-level logical behavior under error correction.

As in \texttt{qstack}, evaluation of \texttt{extern} functions will be treated as opaque: the compiler assumes they return a list of quantum instructions but does not attempt to verify them.  
The output of callbacks will not be verified as part of the compilation process; I expect this is not a blocker since they should typically be used to handle only error-correction feedback, not core program semantics.

\subsection{Compiler Framework}

The compiler will be structured around two categories of passes that differ in both scope and verification method. Local, instruction-level passes will be statically verified once, while global, whole-circuit passes will rely on \emph{translation validation}.

\paragraph{Instruction-level passes.}
These are \emph{local transformations} that operate on one instruction at a time.  
Each pass defines a rule that replaces a specific instruction pattern with an equivalent one.  
For example, a canonicalization pass may replace every occurrence of \(H\) with the sequence \(S\!-\!SX\!-\!S\), leaving all other gates unchanged.  
These transformations are defined syntactically over single gates and are designed to preserve semantics in all contexts.  
Instruction-level passes include canonicalization, gate decomposition, and introduction of encoded (logical) instructions.  
Because each rule acts locally, these passes compose naturally to form larger compilation pipelines.

\paragraph{Whole-circuit passes.}
These are \emph{global transformations} that operate on an entire circuit at once.  
They may restructure or synthesize large regions of code.  
Typical examples include circuit optimization, scheduling, and layout mapping.  
Since these passes transform the full circuit in a single step, they are verified through direct equivalence checking between the input and output circuits rather than by composition of local proofs.

\subsection{Verification Methodology}

The verification methodology follows directly from this division between local and global transformations.
Instruction-level transformations admit static proofs of correctness because their effects are local and context-independent. In contrast, whole-circuit passes may restructure programs globally, making static proofs impractical; these are validated per compilation run using \textbf{translation validation}, which checks that the transformed circuit is equivalent to its input.

\paragraph{Instruction-level passes.}
Each instruction-level pass is verified in two stages:

\begin{enumerate}
  \item \textbf{Local equivalence proof.}  
  For the transformation under consideration (e.g., \(H \rightarrow S\!-\!SX\!-\!S\)), I'll prove that both sides have identical denotations using  \textbf{stabilizer verification} to show both circuits produce the same tableau transformation on the Pauli group.

  \item \textbf{Lifting to global correctness.}  
  Once local equivalence is established, a \textbf{forward-simulation argument} based on the operational semantics is used to prove that applying this transformation to every matching instance throughout a program preserves the program’s overall semantics.  
  Since each step preserves semantics, their composition preserves correctness for the entire compiled circuit.
\end{enumerate}

These passes are verified once, statically, as part of the compiler definition.

\paragraph{Whole-circuit passes.}
For transformations that rewrite or synthesize an entire circuit, verification is performed by direct \textbf{circuit equivalence checking} similarly by applying \textbf{stabilizer verification} when both input and output are Stabilizer circuits, by comparing their induced logical actions on the Pauli group

Because these transformations are verified as a single step, forward simulation is not required.  
Instead, we will use translation validation to check the equivalence of the full input and output circuits. 

\vspace{2em}

The focus of this thesis will be verification of stabilizer circuits. Verification of non-stabilizer circuits is possible by extending this framework using density-matrix or superoperator equivalence checking. However, this component will serve as a stretch goal, exploring how the same verification principles can generalize beyond stabilizer-based reasoning.

s\section{Methodology and Evaluation}

I will evaluate the framework by designing and implementing a verified compiler that transforms quantum programs across multiple abstraction layers.  
The methodology is divided into three main phases: defining instruction sets, building compilers between them, and verifying correctness using the techniques described in the previous section.

\subsection{Instruction Set Design}

The first step will be to design the instruction sets for each abstraction layer, from high-level logical gates to hardware-specific primitives.

\paragraph{High-level circuits.}
At the highest level, I will use a canonical gate set similar to the \textbf{QIR base profile}, including standard gates such as Pauli operations,  \(I\), \(X\), \(Z\), arbitrary rotations, \(RX\), \(RY\), \(RZ\), and basic 2 qubit gates, like \(CX\), and \(CZ\).  
This set will serve as the logical, unencoded representation used by all programs prior to error correction.

\paragraph{Error-correction codes.}
Each quantum error-correcting code defines its own \textbf{instruction set (IS)} that reflects the list of logical primitives it exposes.  
Each IS will be universal, but the exposed primitives differ based on the encoding structure, for example:
\begin{itemize}
    \item \textbf{Steane code \code{7,1,3}:} encodes one logical qubit per patch.  
    Its instruction set will include 1-qubit logical operations (\(X\), \(Z\), \(H\)) and 2-qubit logical operations (\(CX\), \(CZ\)).
    \item \textbf{\code{4,2,2} code}: encodes two logical qubits per patch.  
    It exposes logical two-qubit operations such as \(HH\), \(CX\), \(IX\), and \(ZI\), but does not include single-qubit operations like \(X\) or \(Z\).
    \item \textbf{Surface code}: exposes lattice-surgery primitives such as \emph{split}, \emph{merge}, and \emph{move}, rather than explicit Pauli operations.  
    This reflects the physical operations used to implement logical computation on topological architectures.
\end{itemize}

The process of relating standard gate operations to the corresponding error-corrected primitives is known as \textbf{synthesis}.  
In general, synthesis is non-trivial, but for this project I will implement only trivial examples to illustrate the process.

\paragraph{Intermediate \code{1,1,1} code.}
I will also define an instruction set corresponding to a simple \code{1,1,1} code.  
This code will use a \textbf{Clifford+T} gate set and act as a \emph{generic hardware abstraction layer}.  It is basically a subset of the standard instruction set.
It will primarily serve as an intermediate compilation target when defining circuits that implement logical operations for other QEC codes.

\paragraph{Hardware instruction sets.}
Finally, I will define instruction sets that reflect the physical primitives of real hardware:
\begin{itemize}
    \item \textbf{Ion-trap hardware gates}, including Mølmer–Sørensen (\(R_{ZZ}\)) and local rotation operations;
    \item \textbf{Neutral-atom hardware gates}, based on Rydberg blockade operations and native entangling gates.
\end{itemize}
These instruction sets will serve as final hardware-specific targets in the compilation chain.

\subsection{Compiler Construction}

Once all instruction sets are defined, I will build a family of compilers that map programs between these abstraction layers:
\begin{itemize}
    \item \textbf{Standard $\rightarrow$ QEC primitives:}  
    I will implement proof-of-concept synthesizers for each QEC code that translate high-level gates into their corresponding encoded operations.
    \item \textbf{QEC codes $\rightarrow$ \code{1,1,1}:}  
    I will use the \code{1,1,1} instruction set as a \emph{common intermediate layer} before targeting hardware-specific instruction sets.  
    This ensures a uniform representation for all encoded circuits prior to hardware mapping.
    \item \textbf{\code{1,1,1} $\rightarrow$ Hardware (Neutral or Ion):}  
    I will implement translation passes that adapt the generic Clifford+T representation to the primitives supported by each hardware platform.
\end{itemize}

Each compiler will be implemented as a \textbf{list of verified passes}, where each pass handles a specific primitive.  
This modular design makes it possible to verify each transformation individually and reuse proofs across different compilers.  
For compilers that operate on full circuits (such as \emph{Standard $\rightarrow$ QEC primitives}), I will include automated verification mechanisms that compare the input and output circuits directly to ensure semantic equivalence.

\subsection{Evaluation Plan}

After all compilers are implemented, I will evaluate the framework on multiple representative compilation chains.  
These will include at least:
\begin{itemize}
    \item Standard $\rightarrow$ Steane $\rightarrow$ \code{1,1,1} $\rightarrow$ Ions,
    \item Standard $\rightarrow$ \code{4,2,2} $\rightarrow$ \code{3,1,2} $\rightarrow$  \code{1,1,1} $\rightarrow$ Neutral atoms.
    \item Surface code  $\rightarrow$  \code{1,1,1},
\end{itemize}

The goal of these experiments is to demonstrate that verified passes compose correctly and that the final compiled circuits preserve the semantics of their high-level descriptions across all abstraction layers.

We will not try to test all combinations, though. 
Because compiling to the Surface code instruction set is substantial work (see \cite{litinski2019game}), I will write programs directly using the Surface code primitives (lattice surgery operations) for these experiments.


\subsection{Faulty Passes}

I will implement a small set of intentionally faulty compiler passes to confirm that the verification machinery detects errors. At minimum, I will include:

\begin{itemize}
    \item An instruction-level pass that \emph{generates an incorrect decomposition};
    \item A syndrome-extraction routine with an \emph{incorrect circuit};
    \item A logical-instruction encoder with an \emph{incorrect implementation};
    \item An optimization pass that \emph{incorrectly removes gates}.
\end{itemize}




\subsection{Scope and Limitations}

I will focus exclusively on the correctness of circuit-level compilation and encoded transformations.  
Classical constructs such as functions, loops, or higher-level control structures will not be included in this initial implementation---these can be added later as surface meta-language without changing the core semantics.
I will also not implement full compiler passes for scheduling or optimization, as these topics fall outside the scope of this work, though the same verification framework should support them in future extensions.

As mentioned earlier, the focus will be verification of stabilizer circuits. Verification of non-stabilizer circuits will serve as a stretch goal.

The final outcome will be a verified compiler that demonstrates correctness across multiple abstraction layers and target hardware models, establishing a foundation for scalable and compositional verification of quantum program compilation.


\bibliographystyle{plain}
\bibliography{references}

\end{document}
